\documentclass{article}

\usepackage{amsmath,amssymb}
\usepackage{xurl}
\usepackage[hidelinks]{hyperref}
\setlength{\emergencystretch}{2em}

\input{command}

\title{The marginal-support expansion in CPSV16 Appendix D.2}
\author{}
\date{2026-07-15}

\begin{document}
\maketitle
\gapnote{local-correction}{resolved}

Let $\mathcal H_A$, $\mathcal H_X$, and $\mathcal H_B$ be finite-dimensional
Hilbert spaces. Let
$K_{AXB}\subseteq\mathcal H_A\otimes\mathcal H_X\otimes\mathcal H_B$, and let
$P_{AXB}$ be the orthogonal projector onto $K_{AXB}$. Define the reduced
operators by
\begin{align}
    \rho_{AX}
        &= \tr_B(P_{AXB}),
    \label{eq:cpsvdecor_reduced_operator_ax}\\
    \rho_{XB}
        &= \tr_A(P_{AXB}).
    \label{eq:cpsvdecor_reduced_operator_xb}
\end{align}
The two marginal supports are $\operatorname{ran}(\rho_{AX})$ and
$\operatorname{ran}(\rho_{XB})$, with orthogonal projectors $P_{AX}$ and
$P_{XB}$, respectively. For an operator on $AX$ or $XB$, a hat denotes its
lift to the full tripartite space. In particular,
\begin{align}
    \widehat P_{AX}
        &= P_{AX}\otimes\Id_B,
    \label{eq:cpsvdecor_lifted_projector_ax}\\
    \widehat P_{XB}
        &= \Id_A\otimes P_{XB}.
    \label{eq:cpsvdecor_lifted_projector_xb}
\end{align}
The same convention defines $\widehat\rho_{AX}$ and
$\widehat\rho_{XB}$.

\section{The source step}

In the proof of Proposition~D.3 of
Ref.~\cite[lines 2236--2242]{Cirac2016MPDO_arXiv}, the source chooses
orthonormal bases $\{\alpha_i\}_i$ and $\{\beta_i\}_i$ of the two marginal
supports and asserts that each basis vector is a single partial contraction:
\begin{align}
    \alpha_i
        &= \bra{\widetilde b_i}\varphi_{b,i},
    \label{eq:cpsvdecor_source_single_contraction_ax}\\
    \beta_i
        &= \bra{\widetilde a_i}\varphi_{a,i},
    \label{eq:cpsvdecor_source_single_contraction_xb}
\end{align}
where $\varphi_{a,i},\varphi_{b,i}\in K_{AXB}$ and the contractions act on the
corresponding outer tensor factors.

The support of a reduced operator is spanned by partial contractions, but an
arbitrary vector in that span need not equal one partial contraction. Thus
Eqs.~\ref{eq:cpsvdecor_source_single_contraction_ax} and
\ref{eq:cpsvdecor_source_single_contraction_xb} do not follow from the
definition of marginal support.

\section{Local correction}

The subsequent matrix-unit argument needs only a factorization of each
marginal support projector through the corresponding reduced operator. Fix
orthonormal bases of $\mathcal H_A$ and $\mathcal H_B$, and define the
elementary matrix units on the outer factors by
\begin{align}
    E^A_{ak}
        &= \ketbra{a}{k},
    \label{eq:cpsvdecor_matrix_unit_a}\\
    E^B_{bk}
        &= \ketbra{b}{k}.
    \label{eq:cpsvdecor_matrix_unit_b}
\end{align}
When these matrix units occur in products with $P_{AXB}$, they are lifted by
the identity on the other tensor factors. The reduced operators then have the
exact expansions
\begin{align}
    \widehat\rho_{AX}
        &= \sum_{b,k} E^B_{bk}P_{AXB}E^B_{kb},
    \label{eq:cpsvdecor_reduced_expansion_ax}\\
    \widehat\rho_{XB}
        &= \sum_{a,k} E^A_{ak}P_{AXB}E^A_{ka}.
    \label{eq:cpsvdecor_reduced_expansion_xb}
\end{align}

Define $g(0)=0$ and $g(x)=x^{-1}$ for $x\neq 0$. Functional calculus gives
\begin{align}
    P_{AX}
        &= g(\rho_{AX})\rho_{AX}
         = \rho_{AX}g(\rho_{AX}),
    \label{eq:cpsvdecor_support_factorization_ax}\\
    P_{XB}
        &= g(\rho_{XB})\rho_{XB}
         = \rho_{XB}g(\rho_{XB}).
    \label{eq:cpsvdecor_support_factorization_xb}
\end{align}
The decorrelation identity first gives
\begin{align}
    \widehat\rho_{XB}(\Id-P_{AXB})\widehat\rho_{AX}
        &= 0,
    \label{eq:cpsvdecor_reduced_complement_product_xbax}\\
    \widehat\rho_{AX}(\Id-P_{AXB})\widehat\rho_{XB}
        &= 0.
    \label{eq:cpsvdecor_reduced_complement_product_axxb}
\end{align}
Multiplication by the reciprocal-on-support factors from
Eqs.~\ref{eq:cpsvdecor_support_factorization_ax} and
\ref{eq:cpsvdecor_support_factorization_xb} yields
\begin{align}
    \widehat P_{XB}(\Id-P_{AXB})\widehat P_{AX}
        &= 0,
    \label{eq:cpsvdecor_support_complement_product_xbax}\\
    \widehat P_{AX}(\Id-P_{AXB})\widehat P_{XB}
        &= 0.
    \label{eq:cpsvdecor_support_complement_product_axxb}
\end{align}
The absorption identities then imply
\begin{align}
    \widehat P_{XB}\widehat P_{AX}
        &= P_{AXB},
    \label{eq:cpsvdecor_support_product_xbax}\\
    \widehat P_{AX}\widehat P_{XB}
        &= P_{AXB}.
    \label{eq:cpsvdecor_support_product_axxb}
\end{align}

This correction retains the matrix-unit and marginal-support argument of
Ref.~\cite{Cirac2016MPDO_arXiv}, but replaces the unsupported
single-contraction assertion by the valid factorizations in
Eqs.~\ref{eq:cpsvdecor_support_factorization_ax} and
\ref{eq:cpsvdecor_support_factorization_xb}. It adds no hypothesis to
Proposition~D.3 of Ref.~\cite{Cirac2016MPDO_arXiv}.

\section{Formal statement}

The reduced-operator expansions in
Eqs.~\ref{eq:cpsvdecor_reduced_expansion_ax} and
\ref{eq:cpsvdecor_reduced_expansion_xb} are formalized by
\leanid{TripartiteDecorrelation.lift_marginalAX_eq_sum} and
\leanid{TripartiteDecorrelation.lift_marginalXB_eq_sum}. The complementary
support products in
Eqs.~\ref{eq:cpsvdecor_support_complement_product_xbax} and
\ref{eq:cpsvdecor_support_complement_product_axxb} are formalized by
\leanid{TripartiteDecorrelation.IsDecorrelated.}
\hspace{0pt}\leanid{supports_mul_complement_mul_eq_zero} and its reverse-order
counterpart. The product identities in
Eqs.~\ref{eq:cpsvdecor_support_product_xbax} and
\ref{eq:cpsvdecor_support_product_axxb} are assembled by
\leanid{TripartiteDecorrelation.supportProducts_eq}.

\bibliographystyle{alpha}
\bibliography{references}

\end{document}
